% Part: normal-modal-logic
% Chapter: syntax-and-semantics
% Section: normal-modal-logics

\documentclass[../../../include/open-logic-section]{subfiles}

\begin{document}

\olfileid[pt]{nml}{syn}{nor}

\olsection{Lógicas Modais Normais}

Uma lógica modal pode ser vista como qualquer conjunto adequadamente bem-comportado de
!!{formula}s. Há algumas lógicas modais de particular
interesse, por terem interpretações e
aplicações interessantes, mas principalmente por serem teoricamente bem
compreendidas. Essas são as chamadas lógicas modais normais.

\begin{defn}\ollabel{defn:modal-logic}
Uma \emph{lógica modal}~$\Log L$ é um conjunto de !!{formula}s da linguagem
modal básica que 
\begin{enumerate}
\item contém todas as tautologias proposicionais, 
\item é fechado sob substituição uniforme e 
\item é fechado sob modus ponens, isto é, sempre que $!A$ e $(!A \lif
  !B) \in \Log L$, então também $!B \in \Log L$.
\end{enumerate}
\end{defn}

\begin{ex}
Alguns conjuntos de !!{formula}s reconhecidamente pouco interessantes que
satisfazem essa definição e, portanto, contam como lógicas modais são:
\begin{enumerate}
\item o conjunto de todas as tautologias
\item o conjunto de todas as !!{formula}s (a lógica modal inconsistente)
\end{enumerate}
\end{ex}

\begin{defn}\label{defn:normal-modal-logic}
Uma lógica modal~$\Log L$ é \emph{normal} se, e somente se, 
\begin{enumerate}
\item contém ambas as fórmulas
\begin{align}
\tag{K} \Box(p \lif q) & \lif (\Box p \lif \Box q) \\
\tag{Dual} \Diamond p & \liff \lnot \Box \lnot p
\end{align}
\item e é fechada sob necessitação, isto é, sempre que $!A \in
  \Log L$, então também $\Box !A \in \Log L$.
\end{enumerate}
\end{defn}

Encontraremos duas maneiras principais de definir lógicas modais. Uma é
em termos de teoria da prova, como o conjunto de !!{formula}s que podem ser derivadas em
certos sistemas de !!{derivation}. A outra é em termos de teoria dos modelos, como
conjuntos de !!{formula}s válidas em certas classes de
modelos. Isso reflete a situação na lógica proposicional usual (ou, aliás, na
lógica de primeira ordem). Também aqui podemos determinar um
conjunto de !!{formula}s de duas maneiras diferentes, por exemplo, como o conjunto de
tautologias verdadeiras em todas as valorações e como o
conjunto de todas as !!{derivable} !!{formula}s em algum sistema de !!{derivation}.
No caso das lógicas modais normais, essa dupla maneira de especificar lógicas modais
é possível por meio de modelos relacionais e sistemas de prova
axiomáticos (do tipo Hilbert).
\end{document}
